\section{Sviluppo del modello}\label{sec:sviluppo_del_modello} 


la stesura del codice si è divisa in due momenti fondamentali: la simulazione di segnali (e l'analisi dei suddetti)
e lo studio di immagini reali, acquisite tramite tecniche di microscopia a fluorescenza.\\

Sono state utilizzate le stesse funzioni per lo studio delle due tipologie di segnale.\\





\subsection{Generazione dei segnali}

% vedrò in seguito se spostare questa roba nella parte teorica

Come citato nella sezione (), %aggiungere riferimento ai cenni teorici nella sezione adeguata
i decadienti di fluorescenza seguono un andamento esponenziale decrescente;

\begin{equation}
    N(t) = N_0 e^{-\frac{t}{\tau}}
    \label{eq:exp_decay}
\end{equation}

dove \tau indica il tempo di vita di fluorescenza.\\

Si è dunque cercato di simulare un esperimento di fluorescenza così come avverrebbe in laboratorio: un fascio laser
periodico colpisce il campione causando un'eccitazione elettronica; al tempo \textit{t} un atomo ha una certa 
probabilità di trovarsi ancora allo stato eccitato. La dipendenza temporale di questa probabilità sarà 
dettata da \ref{eq:exp_decay}.\\

% Supponiamo ora di analizzare in tempo reale l'emissione di fluorescenza del campione per mezzo di un microscopio
% a fluorescenza. I fotoni emessi dal campione possono essere raggruppati in \textit{pixel}: l'intensità di ogni
% pixel avrà lo stesso andamento temporale di \ref{eq:exp_decay}.\\

Per ogni pulsazione del laser, si può misurare a quale \textit{t} dal momento dell'impulso avviene la 
diseccitazione (emissione di fluorescenza), e campionare i risultati ottenuti in \textit{t} di larghezza fissata
dalla scheda di acquisizione. Dopo un numero sufficiente di eventi, si sarà ottenuta una distribuzione esponenziale
decrescente, con tempo caratteristico \tau pari al tempo di vita di fluorescenza associato al fluoroforo.\\

Il fenomeno è stato simulato con una semplice generazione di distribuzione esponenziale, e il valore di \tau
verificato mediante il metodo della \textit{Binned Log Likelihood} (è stata utilizzata la libreria \textit{iMinuit}).

% reference teoria iminuit





\subsection{Deconvoluzione del segnale IRF}